\chapter{Two derivations worth doing in full}\label{app:derivations}
\section{Deriving the Zernike radial formula and orthogonality}
Chapter \ref{ch:zernike} constructed low-order modes by subtracting
overlaps. Here is a general construction that proves the normalization
used throughout the book.

Take $m\geq0$, $k\geq0$, $n=m+2k$, and $t=\rho^2$.
Define
\begin{equation}
Q_k(t)=\frac{t^{-m}}{k!}
\frac{\dd^k}{\dd t^k}\left[t^{m+k}(t-1)^k\right],
\qquad R_{m+2k}^m(\rho)=t^{m/2}Q_k(t).
\label{eq:rodrigues}
\end{equation}
Although $t^{-m}$ appears, the derivative contains a factor $t^m$,
so $Q_k$ is an ordinary polynomial of degree $k$.
At $t=1$, only the term in which all $k$ derivatives act on
$(t-1)^k$ survives. Its value is $k!$, giving $Q_k(1)=1$
and $R_n^m(1)=1$.

Expand the bracket:
\[
t^{m+k}(t-1)^k
=\sum_{s=0}^k(-1)^{k-s}\binom ks t^{m+k+s}.
\]
Taking $k$ derivatives gives
\[
R_{m+2k}^m(\rho)
=\sum_{s=0}^k
\frac{(-1)^{k-s}(m+k+s)!}
{s!(k-s)!(m+s)!}\rho^{m+2s}.
\]
Set $j=k-s$, so that $m+2s=n-2j$ and $m+k+s=n-j$.
This becomes
\[
R_n^m(\rho)=\sum_{j=0}^{(n-m)/2}
\frac{(-1)^j(n-j)!}
{j![(n-m)/2-j]![(n+m)/2-j]!}\rho^{n-2j},
\]
which is Eq.~\eqref{eq:radial}.
The factorial formula is therefore the expanded form of a weighted
polynomial orthogonalization construction.

To prove orthogonality, take $\ell<k$ and evaluate
\begin{align}
\int_0^1 R_{m+2k}^mR_{m+2\ell}^m\,\rho\dd\rho
&=\frac12\int_0^1t^mQ_k(t)Q_\ell(t)\dd t\\
&=\frac1{2k!}\int_0^1Q_\ell(t)
\frac{\dd^k}{\dd t^k}[t^{m+k}(t-1)^k]\dd t.
\end{align}
Integrate by parts $k$ times. Boundary terms vanish: the bracket has
sufficient zeros at both $t=0$ and $t=1$.
The remaining integral contains $\dd^kQ_\ell/\dd t^k$, which is zero
because $Q_\ell$ has degree less than $k$.
This proves orthogonality for different radial degrees at fixed $m$.

For the norm, the leading coefficient of $Q_k$ is
\[
A_k=\frac{(m+2k)!}{k!(m+k)!}.
\]
Hence $\dd^kQ_k/\dd t^k=k!A_k$, and integration by parts gives
\begin{align}
\int_0^1[R_{m+2k}^m]^2\rho\dd\rho
&=\frac{A_k}{2}\int_0^1t^{m+k}(1-t)^k\dd t\\
&=\frac{A_k}{2}\frac{(m+k)!\,k!}{(m+2k+1)!}\\
&=\frac{1}{2(m+2k+1)}=\frac1{2(n+1)}.
\end{align}
The middle integral follows by repeated integration by parts:
$\int_0^1t^a(1-t)^b\dd t=a!b!/(a+b+1)!$ for nonnegative integer
$a,b$. Angular sine/cosine orthogonality supplies the remaining
normalization factors in Eq.~\eqref{eq:Zdef}.

\section{Why phase can move edges without changing total energy}
For $U=\FT^{-1}[\widetilde tP]$, Parseval's identity gives
\begin{equation}
\int|U(\vct x)|^2\dd^2x
=\int|\widetilde t(\vct f)|^2|P(\vct f)|^2\dd^2f.
\end{equation}
If only phase changes and $|P|$ stays fixed, the integrated image energy
is unchanged. For oblique illumination replace $P(\vct f)$ by
$P(\vct f+\vct s)$; the same conclusion holds for a fixed source point.
Source averaging preserves it.

Yet local intensity contains cross terms between frequencies.
Their phase differences change with $W$. Energy is redistributed,
altering maxima, minima, slopes, and threshold edges.
Phase aberration does not have to absorb light to degrade printing.

For periodic images, use energy per unit cell rather than an infinite
integral over repeated cells:
\[
\frac1{p_xp_y}\int_{\rm cell}|U|^2\dd x\dd y
=\sum_{m,n}|a_{mn}P_{mn}|^2.
\]
Distinct orders average to zero in the integrated cross terms,
but those cross terms remain essential locally.
If amplitude, aperture support, or coating reflectance changes too,
the total collected energy can change. The phase-only pupil model
isolates one mechanism rather than describing every physical perturbation.

\chapter{Exercises with worked solutions}\label{app:exercises}
Try each problem before reading its solution. Numerical examples use
this book's conventions unless explicitly stated otherwise.

\section{Phase, geometry, and diffraction}
\subsection*{1. Convert milliwaves into length and phase}
\textbf{Problem.} At 13.5 nm, convert $20\,\mathrm{m}\lambda$
into OPD and phase. Does this number alone specify RMS or PV?

\textbf{Solution.}
$W=13.5(20/1000)=0.270\nm$ and
$\phi=2\pi(0.020)=0.125664$ rad.
The conversion gives the size of the stated OPD quantity. Whether it
is a coefficient, RMS, PV, or local error must be specified separately.
For an RMS-normalized coefficient, its magnitude is its RMS contribution.

\subsection*{2. Convert a mirror-height error}
\textbf{Problem.} A local mirror height error is 0.10 nm at $6^\circ$
incidence from the normal. Estimate OPD and phase magnitude.

\textbf{Solution.}
$|\delta W|\simeq2(0.10)\cos6^\circ=0.19890\nm$.
At 13.5 nm this is about 14.734 milliwaves and 0.09257 rad.
The local geometric estimate excludes coating phase and changed ray
intersections.

\subsection*{3. Separate resolution and focus scaling}
\textbf{Problem.} At fixed wavelength, NA increases from 0.33 to 0.55.
How do $\lambda/\NA$ and $\lambda/\NA^2$ change?

\textbf{Solution.}
Their ratios are $0.33/0.55=0.60$ and
$(0.33/0.55)^2=0.36$.
The lateral diffraction scale decreases by 40 percent; the focus
scale decreases by 64 percent. Process-specific constants may also
change, so these are scaling comparisons.

\subsection*{4. Explain why one admitted order cannot draw a contact}
\textbf{Problem.} Show that one transmitted diffraction order has uniform
intensity. Why does the 38 nm example benefit from oblique light?

\textbf{Solution.}
With $U=a\ee^{\ii2\pi(f_xx+f_yy)}$, $|U|^2=|a|^2$ everywhere.
At the stated wavelength and NA, first-order normalized spacing for
38 nm pitch is 1.07656, outside the unit pupil for on-axis illumination.
Oblique source points shift several orders into the pupil, allowing
spatially varying interference cross terms.

\subsection*{5. Recover the phase-to-shift relation}
\textbf{Problem.} Two image orders separated by frequency $1/(38\nm)$
receive relative phase 0.10 rad. Find the fringe shift.

\textbf{Solution.}
$\Delta x=-p\Delta\phi/(2\pi)
=-38(0.10)/(2\pi)\simeq-0.60479\nm$.
This is a differential-phase result for two orders, not a general
coefficient-to-PS calibration.

\subsection*{6. Add amplitudes before intensities}
\textbf{Problem.} Two coherent equal-amplitude fields each have intensity
one. Find total intensity for relative phases 0, $\pi/2$, and $\pi$.
What if the sources are mutually incoherent?

\textbf{Solution.}
$I=2+2\cos\Delta\phi$, giving 4, 2, and 0.
For incoherent contributions the averaged cosine term vanishes and
intensity is 2. Averaging intensities prematurely loses interference.

\section{Wavefront statistics and Zernike modes}
\subsection*{7. Normalize defocus}
\textbf{Problem.} For $W=B\rho^2$, find its mean and piston-removed
RMS. Express the nonpiston part using $Z_2^0$.

\textbf{Solution.}
$\avg W=B/2$ and $\sigma_W^2=B^2(1/3-1/4)=B^2/12$.
The residual is
$B(\rho^2-1/2)=[B/(2\sqrt3)]Z_2^0$.
Its unit-RMS coefficient is $B/(2\sqrt3)$.

\subsection*{8. Distinguish raw and balanced spherical aberration}
\textbf{Problem.} Why is $B\rho^4$ not pure $Z_4^0$?
Find the best quadratic compensation in RMS.

\textbf{Solution.}
The raw quartic overlaps piston and defocus. For
$B\rho^4+d\rho^2$, minimizing piston-removed variance gives $d=-B$.
The residual after mean removal is
$B(\rho^4-\rho^2+1/6)=[B/(6\sqrt5)]Z_4^0$.
Its RMS is $|B|/\sqrt{180}$.

\subsection*{9. Count modes and total RMS}
\textbf{Problem.} How many real disk modes exist through degree six?
What RMS results if 24 orthonormal nonpiston coefficients are all
20 milliwaves?

\textbf{Solution.}
There are $\sum_{n=0}^6(n+1)=28$ modes including piston.
RMS is $\sqrt{24}(20)=97.9796$ milliwaves.
Neither calculation determines the paper's unspecified single-index ordering.

\subsection*{10. Find a pure-tilt placement allowance}
\textbf{Problem.} For $W=cZ_1^{-1}$ and NA 0.33, what $|c|$
corresponds to 0.4 nm vertical shift?

\textbf{Solution.}
$Z_1^{-1}=2v$ gives $|\Delta y|=2|c|/\NA$.
Thus $|c|=0.4(0.33)/2=0.066\nm=4.88889$ milliwaves at 13.5 nm.
This assumes the declared basis and no later removal of the translation.

\subsection*{11. Decide which error a Strehl calculation omits}
\textbf{Problem.} A report removes piston, tilt, and defocus before giving
a high Strehl estimate. Does it guarantee small absolute pattern shift
at a fixed wafer plane?

\textbf{Solution.}
No. Tilt changes absolute placement while preserving registered shape;
defocus compensation changes the observation plane.
The residual quality excludes those freedoms. Placement and fixed-plane
focus must be constrained separately.

\section{Image edges and process variation}
\subsection*{12. Distinguish CD change from PS change}
\textbf{Problem.} Left and right edges move by $-0.2$ and $+0.3$ nm.
Find width and centre changes.

\textbf{Solution.}
$\delta\mathrm{CD}=0.3-(-0.2)=0.5\nm$.
$\delta\mathrm{PS}=[0.3+(-0.2)]/2=0.05\nm$.
Most of the change is widening, with a small net translation.

\subsection*{13. Use the edge-slope formula}
\textbf{Problem.} At a bright feature's right edge, $I'=-0.1$ intensity
units per nm and $\delta I=+0.02$. Find edge displacement.

\textbf{Solution.}
$\delta x=-\delta I/I'=+0.2\nm$.
The edge moves outward. If the left slope is $+0.1$ with the same
perturbation, that edge moves $-0.2\nm$, so width increases 0.4 nm
and the centre stays fixed.

\subsection*{14. Check NILS with an exactly soluble image}
\textbf{Problem.} Let $I(x)=I_0\exp[-x^2/(2s^2)]$ and $T=I_0/2$.
Find nominal width, NILS, and width sensitivity to $\ln D$.

\textbf{Solution.}
The positive edge is $r=s\sqrt{2\ln2}$, width is $w=2r$, and
$|\partial\ln I/\partial x|=r/s^2$.
Thus $\mathrm{NILS}=2r^2/s^2=4\ln2$.
At dose $D$,
\[
w(D)=2s\sqrt{2\ln(2D)},\qquad
\frac{\dd w}{\dd\ln D}=\frac{2s}{\sqrt{2\ln(2D)}}.
\]
At $D=1$, this equals $2w/\mathrm{NILS}$, verifying the two-edge
formula.

\subsection*{15. Show that PVB is not always CD range}
\textbf{Problem.} A 14 nm feature translates between centres
$-0.3$ and $+0.3$ nm over process conditions. Find edge-band widths
and CD range.

\textbf{Solution.}
Each edge translates over 0.6 nm, so $B_L=B_R=0.6\nm$.
Every width remains 14 nm, so $B_{\rm CD}=0$.
Contour variation cannot be inferred from CD range alone.

\section{Data and neural-network reasoning}
\subsection*{16. Audit the sample count}
\textbf{Problem.} Why does the single-mode sweep have 984 rows but at
most 961 distinct wavefronts?

\textbf{Solution.}
$41(24)=984$ rows, but the all-zero vector appears 24 times.
Removing 23 extra copies leaves at most 961 distinct vectors.
Other groups can contain further duplicates.

\subsection*{17. Check normalization and extrapolation}
\textbf{Problem.} A feature trained over $[10,30]$ nm maps to $[-1,1]$.
Find normalized values at 20 and 35 nm.

\textbf{Solution.}
$\tilde x=2(x-10)/20-1$ gives 0 and 1.5.
The latter input is outside the observed range. Scaling does not clip
it or make extrapolation reliable.

\subsection*{18. Count a network's parameters}
\textbf{Problem.} Count parameters in a dense $84\to64\to25$ network
with biases.

\textbf{Solution.}
There are $(84+1)64=5440$ in the first layer and
$(64+1)25=1625$ in the second, totaling 7065.
The paper's longer literal architecture has 24,485.
Input entries are not themselves trainable parameters.

\subsection*{19. Calculate one LM step}
\textbf{Problem.} For $e(\theta)=2\theta-1$, start at $\theta=0$
with damping $\mu=1$. Find the LM step.

\textbf{Solution.}
$e=-1$, $J_e=2$, and
$(2^2+1)\Delta\theta=-2(-1)$.
Hence $\Delta\theta=0.4$ and the new residual is $-0.2$.
With $\mu=0$, the linear problem is solved exactly by $\Delta\theta=0.5$.

\subsection*{20. Explain a valid coefficient mismatch}
\textbf{Problem.} For $F(c_1,c_2)=c_1+c_2$, an original pair is
$(1,0)$ and the inverse proposes $(0,1)$. Is it necessarily wrong?

\textbf{Solution.}
Both produce output 1. Coefficient error is nonzero, but forward
consistency error is zero. An additional indicator such as $c_1-c_2$
would distinguish them. Equivalence depends on the selected requirements.

\subsection*{21. Explain an invalid mean of valid inverses}
\textbf{Problem.} For $F(c)=c^2$, targets $+1$ and $-1$ both yield
input $y=1$. What single prediction minimizes their average squared
coefficient loss, and is it feasible?

\textbf{Solution.}
Loss is $\hat c^2+1$, minimized at $\hat c=0$.
But $F(0)=0$, not 1. Inverse ambiguity can undermine regression even
when every training pair is physically valid.

\section{Validation and budgeting}
\subsection*{22. Compare percentage and absolute error}
\textbf{Problem.} Reference shift is 0.005 nm and prediction is 0.015 nm.
Find absolute and percentage error. Does the prediction meet a 0.4 nm
magnitude limit?

\textbf{Solution.}
Absolute error is 0.010 nm; percentage error is
$100(0.010/0.005)=200\%$.
The predicted 0.015 nm magnitude is below 0.4 nm.
A large percentage need not imply a physical limit violation.

\subsection*{23. Derive a safe two-coefficient box}
\textbf{Problem.} A local shift change is
$\delta y=2\delta c_1-\delta c_2$ and remaining allowance is 0.3.
Assign equal independent bounds $|\delta c_j|\leq a$.

\textbf{Solution.}
Maximum absolute change is $2a+a=3a$, so $a\leq0.1$.
Testing the first coefficient alone would allow 0.15; testing the
second alone would allow 0.3. Those individual limits cannot be
simultaneously permitted independently.

\subsection*{24. Propagate correlated uncertainty}
\textbf{Problem.} Use gradient $\vct g=(1,2)^T$ and covariance
$C=\begin{pmatrix}0.01&0.005\\0.005&0.04\end{pmatrix}\mathrm{nm}^2$
to find metric variance.

\textbf{Solution.}
$\vct g^TC\vct g=0.01+4(0.04)+4(0.005)=0.19\,\mathrm{nm}^2$.
Standard deviation is $0.43589\nm$.
Ignoring covariance gives $0.17\,\mathrm{nm}^2$ and underestimates
variance for this gradient.

\subsection*{25. Identify what remains unknown}
\textbf{Problem.} What essential information is missing for reproducing
every bar in the paper's Fig.~9 exactly?

\textbf{Solution.}
The coefficient convention; mask layouts resolving Tables 2 and 3;
source, stack, polarization, and simulator settings; metric extraction
definitions; training and test arrays; filtering and scaling settings;
network activations, weights, and split; and all 118 target and
verification vectors. The paper explains a workflow but does not
supply these artifacts. A reproduction must obtain or explicitly replace them.

\chapter{Reading map, notation, and source audit}\label{app:reading}
\section{Verified entry points into the supplied textbooks}
Printed page numbers refer to numbers on book pages.
PDF page numbers are one-based viewer positions in the supplied files.
The table gives verified entry points, not an assumed constant offset.

\begin{longtable}{@{}L{.27\textwidth}L{.44\textwidth}L{.20\textwidth}@{}}
\toprule Topic & Section and printed page & PDF page\\ \midrule\endhead
Scalar waves & Born--Wolf \S1.3, p.~14 &43\\
Reflection/refraction & Born--Wolf \S1.5, p.~36 &65\\
Layered media & Born--Wolf \S1.6, p.~51 &80\\
Characteristic matrices & Born--Wolf \S1.6.2, p.~55 &84\\
Periodic layers & Born--Wolf \S1.6.5, p.~66 &95\\
Geometrical-optics limit & Born--Wolf \S3.1, p.~109 &138\\
Stops and pupils & Born--Wolf \S4.8.2, p.~186 &215\\
Wave/ray aberrations & Born--Wolf \S5.1, p.~203 &232\\
Diffraction integral & Born--Wolf \S8.3.1, p.~375 &413\\
Circular-aperture diffraction & Born--Wolf \S8.5.2,
formula on p.~396 &434\\
Microscope resolution & Born--Wolf \S8.6.3, p.~419 &459\\
Aberration diffraction & Born--Wolf \S9.1.1, p.~460 &559\\
Reference-sphere changes & Born--Wolf \S9.1.2, p.~462 &557\\
Small-error intensity & Born--Wolf \S9.1.3, p.~463 &556\\
Zernike definition & Born--Wolf \S9.2.1, p.~464 &555\\
Radial-polynomial table & Born--Wolf \S9.2.1, p.~465 &554\\
Wavefront expansion & Born--Wolf \S9.2.2, p.~466 &553\\
Aberration tolerances & Born--Wolf \S9.3, p.~468 &551\\
Small-error criterion & Born--Wolf \S9.3, p.~469 &550\\
Coherent image transfer & Born--Wolf \S9.5.1, p.~483 &534\\
Incoherent transfer & Born--Wolf \S9.5.2, p.~486 &531\\
Mutual intensity & Born--Wolf \S10.4.1, p.~507 &510\\
Extended incoherent source & Born--Wolf \S10.4.2, p.~508 &509\\
Hopkins propagation & Born--Wolf \S10.4.2(b), p.~512 &505\\
\midrule
Wave-optics postulates & Saleh--Teich \S2.1, p.~40 &62\\
Complex wave amplitude & Saleh--Teich \S2.2, pp.~41--42 &63--64\\
Interference & Saleh--Teich \S2.5, p.~58 &80\\
Spatial harmonics & Saleh--Teich \S4.1, p.~105 &127\\
Angular spectrum & Saleh--Teich \S4.1, pp.~106--107 &128--129\\
Optical Fourier transform & Saleh--Teich \S4.2, p.~116 &138\\
Diffraction & Saleh--Teich \S4.3, p.~121 &143\\
Airy diffraction & Saleh--Teich \S4.3, p.~123 &145\\
Image formation & Saleh--Teich \S4.4, p.~127 &149\\
Coherent transfer function & Saleh--Teich \S4.4, p.~130 &152\\
Defocused pupil & Saleh--Teich \S4.4, p.~133 &155\\
Electromagnetic theory & Saleh--Teich \S5.1, p.~152 &174\\
Reflection/polarization & Saleh--Teich \S6.2, p.~209 &231\\
\bottomrule
\end{longtable}

For a first pass, use Saleh--Teich for waves and Fourier imaging, then
Born--Wolf for wave aberrations, Zernike theory, and partial coherence.
Born--Wolf uses older notation and electromagnetic unit conventions in
places. This guide uses SI for Maxwell's equations and explicitly
declares its complex-phase and polynomial normalization conventions.
Do not equate symbols across books by appearance alone.

The statistical-optics chapter in Saleh--Teich is absent from the
provided file. Its Fourier-transform appendices are also outside the
supplied excerpt. The needed definitions and derivations are included here.

\section{Symbol dictionary}
\begin{longtable}{@{}L{.22\textwidth}L{.53\textwidth}L{.16\textwidth}@{}}
\toprule Symbol & Meaning here & Units\\ \midrule\endhead
$\lambda,\lambda_0$ & Vacuum wavelength; also the wavelength in
vacuum image space &length\\
$k_0$ &$2\pi/\lambda_0$ &inverse length\\
$n,\widetilde n$ &Real or complex refractive index &dimensionless\\
$U,\vct U$ &Scalar or vector complex field amplitude &field scale\\
$I$ &Time-averaged intensity, often normalized &power/area or relative\\
$\mathcal L$ &Optical path length $\int n\dd\ell$ &length\\
$W$ &Actual-minus-reference phase-equivalent OPD &length\\
$\phi$ &Phase $2\pi W/\lambda$ &radians\\
$\sigma_W$ &Piston-removed RMS, unless other removals are stated &length\\
$\rho,\theta$ &Normalized pupil radius and azimuth &dimensionless; rad\\
$u,v$ &Normalized Cartesian pupil coordinates &dimensionless\\
$a$ &Physical pupil radius; opening width in rectangle examples
when stated locally &length\\
$\vct f$ &Transverse spatial-frequency vector &inverse length\\
$f_c$ &Image coherent cutoff $\NA/\lambda$ &inverse length\\
$P$ &Complex coherent pupil transfer function &relative amplitude\\
$S(\vct s)$ &Source weight density &reciprocal coordinate area\\
$Z_n^m$ &This guide's RMS-normalized real basis function &dimensionless\\
$c_n^m$ &Coefficient multiplying that function &length\\
$Z_i,F_i$ &Paper notation: coefficient and basis function &its convention\\
$z$ &Image-space focus displacement &length\\
$D$ &Relative dose multiplier in the contour example &dimensionless\\
$T$ &Fixed exposure threshold &exposure scale\\
$\vct y$ &Vector of imaging indicators &usually nm\\
$F,G$ &Forward imaging map and learned inverse &by arguments\\
$J$ &Imaging Jacobian $\partial\vct y/\partial\vct c$
&metric / coefficient\\
$J_e$ &Training-residual Jacobian versus network parameters &training scale\\
$\Theta$ &All trainable weights and biases &parameter-dependent\\
$\mu$ &LM damping; $\mu_{12}$ denotes coherence elsewhere &context-dependent\\
$B$ &Physical-variable-to-wavefront influence matrix &coefficient / control\\
\bottomrule
\end{longtable}

\section{Publication issues that must remain visible}
\begin{longtable}{@{}L{.31\textwidth}L{.63\textwidth}@{}}
\toprule Issue & Consequence\\ \midrule\endhead
Unspecified Zernike ordering and normalization
&Named high-index modes, RMS conversion, and transfer to another
simulator cannot be fixed uniquely.\\
$Z_2$--$Z_{25}$ versus 25 outputs
&There are 24 listed varied coordinates; the extra output is unexplained.\\
Tables 2 and 3 disagree
&Plot labels cannot be converted into a unique mask prescription.\\
Table 3 first SQ pitch pair
&$(56,56)$ also differs from Table 2's rectangular $(56,70)$.\\
Table 4 repeats PVB\_X
&Interpreting the second as PVB\_Y is plausible but unconfirmed.\\
Table 4 REC horizontal CDs
&Values near 27--30 nm differ strongly from the nominal 14 nm
horizontal size; the cause is unresolved.\\
CRAO and source shaping mixed in prose
&Chief-ray geometry and off-axis source distribution are distinct.\\
Figure 9 extrema exceed training sweep
&Output extrapolation or a scale issue needs clarification;
hard coefficient-box enforcement is not demonstrated.\\
Small-value filtering description
&Retained data, threshold units, and zero treatment are not reproducible.\\
Training split and activations absent
&Layer widths alone do not specify a trained model.\\
Metric extraction not fully defined
&Different CD/PS/PVB conventions can give different bars.\\
Roadmap-to-PS allocation not derived
&The 1 nm motivation is not a universal projector allowance.\\
Data and network weights absent
&The guide reproduces its teaching calculations, not the authors'
exact dataset.\\
\bottomrule
\end{longtable}

\section{A checklist for your own future experiment}
Record wavelength; image and object NA; magnification; physical and
wafer-referred geometry; pupil axes and support; OPD sign and units;
Zernike functions and normalization; removed modes and reference focus;
source weights; polarization; mask stack and materials; image or resist
model; contour measurement rules; dose-focus sampling; coefficient bounds;
training design and duplicate handling; split, scaling, and filtering;
network activations, loss, stopping, and random seed; and forward
verification of each candidate.

That record turns an interesting calculation into something another
researcher can examine, repeat, and improve.
